\section{Concurrency and Data Integrity}

Concurrency protection was moved from Phase 4 to Phase 1 of the V2 roadmap. The reasoning was straightforward: concurrency is a correctness property, not a performance optimisation. A system that may create duplicate bookings under concurrent load is incorrect by definition, regardless of how well it is observed or how cleanly it is deployed. Resolving correctness issues before adding new capabilities avoids building on an unstable foundation.

Two complementary mechanisms were introduced in the Booking Service: optimistic locking to prevent concurrent modification of the same booking, and an idempotency key to prevent duplicate bookings from client retries.

\subsection{Optimistic Locking}

\subsubsection{The Problem}

The booking conflict detection in V1 relied entirely on an application-level overlap check:

\begin{lstlisting}[style=javastyle, language=Java, caption={V1 conflict detection — no concurrency protection}]
boolean hasConflict = repository.existsOverlappingBooking(
    dto.resourceId(), dto.startTime(), dto.endTime()
);
if (hasConflict) {
    throw new BookingConflictException("Resource is already booked during this time.");
}
\end{lstlisting}

Under concurrent load, two requests arriving simultaneously for the same resource and time slot could both execute \texttt{existsOverlappingBooking}, both find no conflict, and both proceed to \texttt{repository.save()} --- creating a double-booking. The k6 load tests in V1 exposed this scenario but did not resolve it.

\subsubsection{Optimistic vs Pessimistic Locking}

Two standard approaches exist for protecting against concurrent modifications: optimistic locking and pessimistic locking.

Pessimistic locking acquires an exclusive lock on a database row at read time using \texttt{SELECT FOR UPDATE}. No other transaction can modify or lock that row until the current transaction commits or rolls back. This guarantees consistency at the cost of throughput: under high concurrency, transactions queue behind the lock holder, degrading performance.

Optimistic locking does not acquire any lock. Instead, it assumes that conflicts are rare and detects them at write time. Each row carries a version counter. When a transaction reads a row, it reads the current version. When it writes, it includes a \texttt{WHERE version = \{read\_version\}} condition. If another transaction has incremented the version in the meantime, the update matches zero rows and the conflict is detected.

The choice between the two depends on the probability and cost of conflicts. Pessimistic locking is appropriate when conflicts are frequent or when the cost of a conflict is high enough to justify the throughput penalty --- financial transactions, inventory with very limited stock, or approval workflows with strict state transitions. Optimistic locking is appropriate when conflicts are rare and retries are acceptable.

For the Booking Service, optimistic locking is the correct choice. Concurrent modifications to the same booking record are unlikely --- the typical scenario is two different clients attempting to book the same resource for the same time slot, which is handled by the overlap check. Optimistic locking adds a safety net without introducing lock contention.

It is worth noting that both mechanisms can coexist in the same system. A future version could introduce pessimistic locking for a high-concurrency offer type --- such as limited-seat events --- while the Booking Service continues to use optimistic locking for standard bookings.

\subsubsection{Implementation}

Optimistic locking in JPA requires a single field annotated with \texttt{@Version}:

\begin{lstlisting}[style=javastyle, language=Java, caption={@Version field added to Booking entity}]
/** optimistic locking — prevents concurrent modification of the same booking */
@Version
private Long version;
\end{lstlisting}

Hibernate manages this field automatically. On every \texttt{UPDATE}, Hibernate appends a version check to the generated SQL:

\begin{lstlisting}[style=sqlstyle, caption={SQL generated by Hibernate for optimistic locking}]
UPDATE bookings
SET status = ?, version = 2
WHERE id = ? AND version = 1
\end{lstlisting}

If the \texttt{WHERE} clause matches zero rows --- because another transaction already incremented the version --- Hibernate throws \texttt{ObjectOptimisticLockingFailureException}. The application never touches the \texttt{version} field directly.

The column was added via a Flyway migration:

\begin{lstlisting}[style=sqlstyle, caption={V2\_\_add\_version\_to\_bookings.sql}]
ALTER TABLE bookings ADD COLUMN version BIGINT NOT NULL DEFAULT 0;
\end{lstlisting}

The \texttt{GlobalExceptionHandler} was updated to handle the exception and return a meaningful HTTP response:

\begin{lstlisting}[style=javastyle, language=Java, caption={Handler for optimistic locking conflicts}]
/**
 * Handles optimistic locking conflicts — thrown when two concurrent requests
 * attempt to modify the same booking simultaneously.
 * Returns HTTP 409 so the client knows to retry.
 */
@ExceptionHandler(ObjectOptimisticLockingFailureException.class)
public ResponseEntity<ApiError> handleOptimisticLocking(
        ObjectOptimisticLockingFailureException ex) {
    ApiError apiError = new ApiError(
            HttpStatus.CONFLICT.value(),
            "Conflict",
            "Booking was modified by another request. Please try again.",
            LocalDateTime.now(),
            Collections.emptyList()
    );
    return ResponseEntity.status(HttpStatus.CONFLICT).body(apiError);
}
\end{lstlisting}

\subsection{Idempotency Key}

\subsubsection{The Problem}

Optimistic locking protects against concurrent modifications of the same record. It does not protect against a different class of problem: client retries after a network failure.

Consider the following sequence of events:

\begin{enumerate}
    \item Client sends \texttt{POST /api/bookings}
    \item Server receives the request, validates it, and creates the booking
    \item Server sends the response
    \item The network fails before the response reaches the client
    \item Client receives a timeout or connection error
    \item Client does not know whether the booking was created
    \item Client retries the request
    \item Server creates a second booking --- a duplicate
\end{enumerate}

From the server's perspective, both requests are valid. From the client's perspective, the retry appears necessary. The result is a duplicate booking that neither side intended.

Without an idempotency key, the client faces a second problem: if it retries and receives HTTP 409 Booking Conflict, it cannot determine whether the conflict is because its own booking already exists or because another client booked the same slot. The error is ambiguous.

\subsubsection{Implementation}

The idempotency key is an optional client-provided string that uniquely identifies a request. If the server receives a request with a key it has already processed, it returns the original response without creating a duplicate.

The key was added to the \texttt{BookingRequestDTO}:

\begin{lstlisting}[style=javastyle, language=Java, caption={Idempotency key added to request DTO}]
public record BookingRequestDTO(
        UUID resourceId,
        String customerName,
        String customerEmail,
        LocalDateTime startTime,
        LocalDateTime endTime,
        BookingStatus status,
        String idempotencyKey // optional — client-provided unique ID to prevent duplicate bookings on retry
) {}
\end{lstlisting}

The \texttt{Booking} entity stores the key with a unique constraint:

\begin{lstlisting}[style=javastyle, language=Java, caption={Idempotency key field in Booking entity}]
/** idempotency key — client-provided unique ID to prevent duplicate bookings on retry */
@Column(unique = true)
private String idempotencyKey;
\end{lstlisting}

The column was added via a Flyway migration:

\begin{lstlisting}[style=sqlstyle, caption={V3\_\_add\_idempotency\_key\_to\_bookings.sql}]
ALTER TABLE bookings ADD COLUMN idempotency_key VARCHAR(255) UNIQUE;
\end{lstlisting}

The \texttt{BookingServiceImpl} checks for an existing booking with the same key before processing the request:

\begin{lstlisting}[style=javastyle, language=Java, caption={Idempotency check in BookingServiceImpl}]
if (dto.idempotencyKey() != null && !dto.idempotencyKey().isBlank()) {
    Optional<Booking> existing =
        repository.findByIdempotencyKey(dto.idempotencyKey());
    if (existing.isPresent()) {
        // return the original response without creating a duplicate
        return mapToResponseDTO(existing.get());
    }
}
\end{lstlisting}

If the key is found, the original booking is returned immediately. The overlap check, the Catalog Service call, and the \texttt{repository.save()} are all bypassed. The client receives the same response it would have received on the first successful request.

\subsection{Defence in Depth}

The two mechanisms are complementary and address different failure modes:

\begin{table}[H]
\centering
\begin{tabular}{lll}
\toprule
\textbf{Mechanism} & \textbf{Protects against} & \textbf{Response} \\
\midrule
Optimistic locking & Concurrent modification of the same booking & HTTP 409 --- retry \\
Idempotency key & Duplicate creation from client retry & HTTP 200 --- original response \\
Overlap check & Double-booking the same resource slot & HTTP 409 --- conflict \\
\bottomrule
\end{tabular}
\caption{Concurrency protection layers in the Booking Service}
\end{table}

Together, they form a defence-in-depth approach to booking integrity. No single mechanism covers all failure modes; each layer addresses a specific scenario.

\subsection{Resource Cache Seeder — Design Decisions}

\subsubsection{Component Classification}

The \texttt{ResourceCacheSeeder} is annotated with \texttt{@Component} rather than \texttt{@RestController}, \texttt{@Service}, or \texttt{@Repository}. Understanding the distinction matters.

Spring's stereotype annotations share a common purpose --- registering a class as a bean in the application context --- but carry different semantic meaning:

\begin{table}[H]
\centering
\begin{tabular}{lp{9cm}}
\toprule
\textbf{Annotation} & \textbf{Intended use} \\
\midrule
\texttt{@RestController} & Exposes HTTP endpoints --- receives requests from external clients \\
\texttt{@Service} & Encapsulates business logic --- called by controllers or other services \\
\texttt{@Repository} & Manages data access --- interacts with the database \\
\texttt{@Component} & Generic bean --- infrastructure or utility components that do not fit the above \\
\bottomrule
\end{tabular}
\caption{Spring stereotype annotations and their intended roles}
\end{table}

The \texttt{ResourceCacheSeeder} does not expose endpoints, does not contain business logic, and does not access the database directly. It is an infrastructure component that orchestrates a startup task. \texttt{@Component} is the correct choice --- honest about what the class is, without overstating its role.

\texttt{@Service} would also be technically valid, but would imply that the seeder contains business logic, which it does not. Choosing the most accurate annotation reduces cognitive load for anyone reading the code.

\subsubsection{ApplicationReadyEvent vs PostConstruct}

The seeder uses \texttt{@EventListener(ApplicationReadyEvent.class)} to trigger the cache population. An alternative would be \texttt{@PostConstruct}, which runs a method immediately after the bean is initialised.

\texttt{@PostConstruct} was not used because it runs too early in the startup sequence --- before Eureka registration is complete and before OpenFeign clients are fully configured. Attempting to call the Catalog Service at this point would fail because the service has not yet been discovered.

\texttt{ApplicationReadyEvent} fires after the entire application context is initialised, all beans are ready, Eureka has registered the service, and the embedded server is accepting connections. At this point, the Feign client can safely resolve and call the Catalog Service.

The distinction is important in any Spring Boot application that makes network calls during startup: \texttt{@PostConstruct} runs during bean initialisation, \texttt{ApplicationReadyEvent} runs after the application is fully started.

\subsection{Fail Fast — Request Processing Order}

The \texttt{BookingServiceImpl.create()} method processes validations in a deliberate order designed to reject invalid requests as early and as cheaply as possible:

\begin{enumerate}
    \item \textbf{Idempotency key check} --- if the request was already processed, return the original response immediately. No database queries beyond the key lookup, no network calls.
    \item \textbf{Overlap check} --- verify that no conflicting booking exists for the requested time slot. Rejected here if a conflict is found, before any call to the Catalog Service.
    \item \textbf{Resource resolution} --- only after confirming the request is new and the time slot is available does the system fetch resource data, either from the local cache or via OpenFeign.
\end{enumerate}

This ordering follows the fail fast principle: reject the request at the earliest possible point, with the lowest possible cost. Inverting the order --- fetching resource data before checking for conflicts --- would result in Catalog Service calls for requests that are ultimately rejected, wasting network resources and adding latency.

The same principle applies to the overlap check preceding the resource resolution: a conflict check is a single database query; resolving a resource may involve a cache lookup, a Feign call, and a cache write. The cheaper operation always comes first.

\subsection{Trade-offs}

Optimistic locking increases the complexity of the update path. Under high concurrency, clients may receive HTTP 409 responses and need to retry. This is acceptable for the Booking Service, where concurrent modifications to the same record are rare. For a system with frequent contention on the same records, pessimistic locking would be more appropriate.

The idempotency key is optional by design. Clients that do not provide a key receive no duplicate protection --- the server cannot distinguish between an intentional second booking and an unintentional retry. Making the key mandatory would provide stronger guarantees but would be a breaking change for existing clients. The current design favours backward compatibility.

\subsection{Physical Delete and Orphaned Bookings}

\subsubsection{The Problem Discovered}

During the implementation of the Kafka hybrid model, a significant design flaw was identified in the original \texttt{DELETE /api/resources/\{id\}} endpoint: it physically deleted a resource from the Catalog Service database without verifying whether active bookings existed in the Booking Service.

This is a direct consequence of the database-per-service pattern. The Booking Service stores only a \texttt{resourceId} reference --- there is no foreign key constraint between the two databases. If a resource is physically deleted from the Catalog, any booking that references it becomes an orphaned record: the booking exists in the Booking Service database, but the resource it references no longer exists anywhere.

From the user's perspective, this means having a confirmed reservation with no associated details --- no service name, no location, no price. This is one of the most damaging user experiences possible in a booking platform.

\subsubsection{The Decision}

Physical delete was removed entirely and replaced with a soft delete pattern. Two new endpoints were introduced:

\begin{itemize}
    \item \texttt{PATCH /api/resources/\{id\}/deactivate} --- marks the resource as inactive, publishes \texttt{RESOURCE\_DEACTIVATED} to Kafka, preserves the record in the database
    \item \texttt{PATCH /api/resources/\{id\}/activate} --- reverses deactivation, publishes \texttt{RESOURCE\_ACTIVATED} to Kafka
\end{itemize}

A deactivated resource remains in the Catalog database indefinitely. Existing bookings retain access to all resource data. New bookings are rejected because the Booking Service cache reflects the inactive state via the Kafka event.

\subsubsection{The Correct Long-Term Solution}

The current approach prevents orphaned bookings but does not support physical delete. The correct solution requires the Booking Service to publish events when bookings complete or are cancelled --- \texttt{BOOKING\_COMPLETED}, \texttt{BOOKING\_CANCELLED} --- and the Catalog Service to maintain a counter of active bookings per resource. Physical delete would only be permitted when this counter reaches zero.

This requires bidirectional Kafka communication between the two services and was deferred to V3, where the system evolves towards a more purely event-driven model.

\subsubsection{Data Retention}

Preserving inactive resources indefinitely has a storage cost. In a production system, this would be governed by a data retention policy --- a scheduled process that archives or purges records after a defined period, conditional on the absence of active bookings. This is a known limitation of the current implementation and a V3 item.

The trade-off is deliberate: additional storage cost is preferable to orphaned bookings. A confirmed reservation with no associated data is a broken user experience; a slightly larger database is an operational concern with straightforward solutions.

\subsection{Lessons Learned}

Concurrency is a correctness property, not a performance concern. The decision to move optimistic locking and the idempotency key from Phase 4 to Phase 1 was correct. Introducing Kafka or Kubernetes on top of a system that could silently create duplicate bookings would have been building on an incorrect foundation.

The idempotency key scenario --- network failure after server success, before client acknowledgement --- is a category of failure that is easy to overlook during development because it does not appear in unit tests. The client either receives a response or it does not. Unit tests always receive a response. The failure only becomes visible under real network conditions, which makes it important to reason about explicitly during design rather than after deployment.